% Generated by hcsm.collect_tensor -- do not edit by hand.
\documentclass[11pt]{article}
\usepackage{amsmath,amssymb,booktabs,graphicx}
\usepackage[margin=1in]{geometry}
\title{Anisotropic tensor diffusion on a torus:\\
which error components defeat point smoothers,\\
and what a learned smoother does about them}
\date{}
\begin{document}
\maketitle

\section{What was run}

The problem is
\[
-\nabla_\Gamma\cdot(D_\varepsilon\nabla_\Gamma u)+u=f_\varepsilon,
\qquad
D_\varepsilon=\varepsilon\,e_\xi\otimes e_\xi+e_\eta\otimes e_\eta,
\]
on the torus $R=2$, $r=1$, with the manufactured solution
$u=\sin\xi\sin\eta$ and
\[
f_\varepsilon=\sin\xi\sin\eta\left(2+\frac{\cos\eta}{h}+\frac{\varepsilon}{h^2}\right),
\qquad h=2+\cos\eta ,
\]
derived in closed form for this tensor and verified before use. The strong
direction is $\eta$ (weight $1$), the weak direction is $\xi$ (weight
$\varepsilon$), so the anisotropy is $1/\varepsilon$. The case
$\varepsilon=1$ is the identity tensor and serves as the isotropic control.


\subsection{Operator verification}

The exported prolongation was compared against \texttt{MGTransferPrebuilt}, the
transfer the V-cycle actually uses, at every level and every $\varepsilon$: the
maximum difference is exactly zero, for both the prolongation and the
restriction. The assembled operator is affine in $\varepsilon$,
$A_\varepsilon=A_\eta+\varepsilon A_\xi$, to a relative $4\times10^{-16}$, which
is what it means for the tensor to be applied consistently; and
$\varepsilon$ does change the assembly, by $33\%$ in max-norm between
$\varepsilon=1$ and $\varepsilon=10^{-4}$. The manufactured solution converges at
$O(h)$ in $H^1$ and $O(h^2)$ in $L^2$ at \emph{every} $\varepsilon$
(table~\ref{tab:conv}).

\begin{table}[htbp]\centering
\caption{Discretisation quality and operator anisotropy. The discretisation is
unaffected by $\varepsilon$; the operator's diagonal spread and coarse-grid
condition number grow only mildly.}
\label{tab:conv}
\begin{tabular}{lrrr}
\toprule
$\varepsilon$ & anisotropy & $\max_i A_{ii}/\min_i A_{ii}$ & $\operatorname{cond}(P^{\mathsf T}A_\varepsilon P)$ \\
\midrule
$1$ & $1{:}1$ & 1.66 & 41.7 \\
$10^{-1}$ & $10{:}1$ & 2.72 & 41.7 \\
$10^{-2}$ & $100{:}1$ & 2.93 & 94.7 \\
$10^{-3}$ & $1000{:}1$ & 2.95 & 121.2 \\
$10^{-4}$ & $10000{:}1$ & 2.96 & 124.7 \\
\bottomrule
\end{tabular}
\end{table}


\section{Classical smoothers: which components are hard}

Norm ratios $\lVert\cdot\rVert_{A_\varepsilon}/\lVert e\rVert_{A_\varepsilon}$,
one application each. ``full'' is the whole error, ``complement'' is
$\lVert Q_\varepsilon e^+\rVert_{A_\varepsilon}/\lVert e\rVert_{A_\varepsilon}$,
the part the coarse grid cannot reach.

\begin{table}[htbp]\centering
\caption{Mean full norm ratio, jacobi, on the held-out test families (sweeps: 1 (damped, omega from validation)).}
\label{tab:cls-jacobi-full}
\begin{tabular}{lrrrrr}
\toprule
error family & $1$ & $10^{-1}$ & $10^{-2}$ & $10^{-3}$ & $10^{-4}$ \\
\midrule
gaussian & 0.393 & 0.517 & 0.536 & 0.538 & 0.538 \\
fourier & 0.379 & 0.467 & 0.473 & 0.473 & 0.473 \\
weak & 0.574 & 0.900 & 0.956 & 0.961 & 0.961 \\
survivor & 0.651 & 0.835 & 0.858 & 0.859 & 0.859 \\
\bottomrule
\end{tabular}
\end{table}

\begin{table}[htbp]\centering
\caption{Mean complement norm ratio, jacobi, on the held-out test families (sweeps: 1 (damped, omega from validation)).}
\label{tab:cls-jacobi-complement}
\begin{tabular}{lrrrrr}
\toprule
error family & $1$ & $10^{-1}$ & $10^{-2}$ & $10^{-3}$ & $10^{-4}$ \\
\midrule
gaussian & 0.363 & 0.471 & 0.487 & 0.488 & 0.488 \\
fourier & 0.350 & 0.423 & 0.426 & 0.425 & 0.425 \\
weak & 0.567 & 0.900 & 0.956 & 0.961 & 0.961 \\
survivor & 0.630 & 0.818 & 0.839 & 0.840 & 0.841 \\
\bottomrule
\end{tabular}
\end{table}

\begin{table}[htbp]\centering
\caption{Mean full norm ratio, gs, on the held-out test families (sweeps: 1 forward sweep).}
\label{tab:cls-gs-full}
\begin{tabular}{lrrrrr}
\toprule
error family & $1$ & $10^{-1}$ & $10^{-2}$ & $10^{-3}$ & $10^{-4}$ \\
\midrule
gaussian & 0.334 & 0.431 & 0.440 & 0.441 & 0.441 \\
fourier & 0.314 & 0.416 & 0.424 & 0.424 & 0.424 \\
weak & 0.388 & 0.786 & 0.904 & 0.915 & 0.916 \\
survivor & 0.502 & 0.673 & 0.688 & 0.689 & 0.689 \\
\bottomrule
\end{tabular}
\end{table}

\begin{table}[htbp]\centering
\caption{Mean complement norm ratio, gs, on the held-out test families (sweeps: 1 forward sweep).}
\label{tab:cls-gs-complement}
\begin{tabular}{lrrrrr}
\toprule
error family & $1$ & $10^{-1}$ & $10^{-2}$ & $10^{-3}$ & $10^{-4}$ \\
\midrule
gaussian & 0.273 & 0.357 & 0.362 & 0.362 & 0.362 \\
fourier & 0.251 & 0.336 & 0.339 & 0.338 & 0.338 \\
weak & 0.361 & 0.785 & 0.903 & 0.915 & 0.916 \\
survivor & 0.455 & 0.646 & 0.661 & 0.661 & 0.661 \\
\bottomrule
\end{tabular}
\end{table}

\begin{table}[htbp]\centering
\caption{Mean full norm ratio, sgs_ssor, on the held-out test families (sweeps: 1 symmetric application = forward + backward sweep).}
\label{tab:cls-sgs_ssor-full}
\begin{tabular}{lrrrrr}
\toprule
error family & $1$ & $10^{-1}$ & $10^{-2}$ & $10^{-3}$ & $10^{-4}$ \\
\midrule
gaussian & 0.163 & 0.241 & 0.250 & 0.251 & 0.251 \\
fourier & 0.149 & 0.232 & 0.241 & 0.241 & 0.241 \\
weak & 0.223 & 0.643 & 0.827 & 0.848 & 0.850 \\
survivor & 0.310 & 0.522 & 0.550 & 0.551 & 0.551 \\
\bottomrule
\end{tabular}
\end{table}

\begin{table}[htbp]\centering
\caption{Mean complement norm ratio, sgs_ssor, on the held-out test families (sweeps: 1 symmetric application = forward + backward sweep).}
\label{tab:cls-sgs_ssor-complement}
\begin{tabular}{lrrrrr}
\toprule
error family & $1$ & $10^{-1}$ & $10^{-2}$ & $10^{-3}$ & $10^{-4}$ \\
\midrule
gaussian & 0.145 & 0.223 & 0.231 & 0.231 & 0.231 \\
fourier & 0.132 & 0.215 & 0.222 & 0.222 & 0.222 \\
weak & 0.202 & 0.641 & 0.827 & 0.848 & 0.850 \\
survivor & 0.287 & 0.513 & 0.540 & 0.541 & 0.541 \\
\bottomrule
\end{tabular}
\end{table}


\section{The learned smoother}

Same errors, same metric. The DeepONet is the framework's residual-to-correction
network with its learned Jacobi skip, wrapped in a scale-equivariant
normalisation so that $B(\lambda r)=\lambda B(r)$ exactly; the measured
equivariance error over amplitudes $10^{-6}$ to $10^{6}$ is below $10^{-15}$.
The loss is
$\mathbb{E}\lVert Q_\varepsilon e^+\rVert_A^2+\lambda\,\mathbb{E}\lVert e^+\rVert_A^2$
with $\lambda$ selected on validation errors.

\begin{table}[htbp]\centering
\caption{Mean full norm ratio, DeepONet, on the held-out test families.}
\label{tab:dn-full}
\begin{tabular}{lrrrrr}
\toprule
error family & $1$ & $10^{-1}$ & $10^{-2}$ & $10^{-3}$ & $10^{-4}$ \\
\midrule
gaussian & 0.413 & 0.441 & 0.438 & 0.432 & 0.425 \\
fourier & 0.401 & 0.469 & 0.471 & 0.475 & 0.462 \\
weak & 0.484 & 0.892 & 0.933 & 0.956 & 0.955 \\
survivor & 0.605 & 0.713 & 0.709 & 0.697 & 0.695 \\
\bottomrule
\end{tabular}
\end{table}

\begin{table}[htbp]\centering
\caption{Mean complement norm ratio, DeepONet, on the held-out test families.}
\label{tab:dn-complement}
\begin{tabular}{lrrrrr}
\toprule
error family & $1$ & $10^{-1}$ & $10^{-2}$ & $10^{-3}$ & $10^{-4}$ \\
\midrule
gaussian & 0.384 & 0.405 & 0.402 & 0.397 & 0.393 \\
fourier & 0.372 & 0.433 & 0.436 & 0.437 & 0.428 \\
weak & 0.464 & 0.874 & 0.925 & 0.953 & 0.951 \\
survivor & 0.580 & 0.697 & 0.689 & 0.681 & 0.680 \\
\bottomrule
\end{tabular}
\end{table}


\section{Held-out stress set}

Analytic near-worst-case modes $\cos(m\xi)\cos(n\eta)$, projected by
$Q_\varepsilon$ and normalised. These are never used for training, model
selection or thresholding.

\begin{table}[htbp]\centering
\caption{Mean full norm ratio on the held-out stress modes, damped Jacobi.}
\label{tab:stress-jacobi}
\begin{tabular}{lrrrrr}
\toprule
error family & $1$ & $10^{-1}$ & $10^{-2}$ & $10^{-3}$ & $10^{-4}$ \\
\midrule
gaussian & -- & -- & -- & -- & -- \\
fourier & -- & -- & -- & -- & -- \\
weak & -- & -- & -- & -- & -- \\
survivor & -- & -- & -- & -- & -- \\
\bottomrule
\end{tabular}
\end{table}

\begin{table}[htbp]\centering
\caption{Mean full norm ratio on the held-out stress modes, DeepONet.}
\label{tab:stress-dn}
\begin{tabular}{lrrrrr}
\toprule
error family & $1$ & $10^{-1}$ & $10^{-2}$ & $10^{-3}$ & $10^{-4}$ \\
\midrule
gaussian & -- & -- & -- & -- & -- \\
fourier & -- & -- & -- & -- & -- \\
weak & -- & -- & -- & -- & -- \\
survivor & -- & -- & -- & -- & -- \\
\bottomrule
\end{tabular}
\end{table}


\section{Actual multigrid performance}

This is a different question from every table above. Here the V-cycle is run on
the exported hierarchy from random right-hand sides, as a stationary defect
correction with no Krylov acceleration, and the contraction rate is measured
from the last cycles. The learned smoother is nonlinear, so where an accelerated
outer solver is wanted the compatible one is flexible CG; ordinary PCG is used
for the classical smoothers.

\begin{table}[htbp]\centering
\caption{V-cycle contraction rates and outer-solve iteration counts on the
1024-DoF level. A rate near $1$ means the cycle barely contracts.}
\label{tab:vcycle}
\begin{tabular}{lrrrrrr}
\toprule
$\varepsilon$ & \multicolumn{3}{c}{stationary rate} & \multicolumn{3}{c}{outer iterations} \\
\cmidrule(lr){2-4}\cmidrule(lr){5-7}
 & Jacobi & SGS & learned & SGS (PCG) & learned (FCG) & Jacobi \\
\midrule
$1$ & 0.3159 & 0.7510 & 0.6974 & 18 & 11 & 28 \\
$10^{-1}$ & 0.7566 & 0.9295 & 0.9403 & 30 & 22 & 500 \\
$10^{-2}$ & 0.8866 & 0.9644 & 0.9708 & 30 & 35 & 500 \\
$10^{-3}$ & 0.9177 & 0.9722 & 0.9776 & 30 & 43 & 500 \\
$10^{-4}$ & 0.9222 & 0.9733 & 0.9784 & 30 & 45 & 124 \\
\bottomrule
\end{tabular}
\end{table}


\section{What is and is not established}

\textbf{Established.} The anisotropic tensor makes the classical point smoothers
fail, and the failure is in the smoother rather than in the discretisation or in
the transfer, both of which are verified independently. The difficult components
are identified: errors oscillatory in $\xi$ and smooth in $\eta$, i.e.
oscillatory along the weak-coupling direction, exactly the regime the standard
anisotropic analysis predicts. The softest evidence is the quantity
$\eta(e)=\lVert D^{-1}Ae\rVert_A/\lVert e\rVert_A$, which enters the bound
$\lVert S_Je\rVert_A\ge(1-\omega\eta(e))\lVert e\rVert_A$; where
$\omega\eta(e)\ge1$ the bound guarantees nothing and damped Jacobi indeed
reduces nothing.

\textbf{Not established.} That the learned smoother is better than the classical
methods overall, or that it rescues the multigrid. The per-family and
per-$\varepsilon$ tables say exactly how much it helps and where; the V-cycle
table says whether that help survives into the solver. No claim is made beyond
this torus, this discretisation, this coefficient family, this training
distribution and this network configuration.

\textbf{Not measured here.} A second mesh: the hierarchy is one fixed level
(1024 DoFs) because the branch network consumes a fixed-length residual and is
bound to one DoF count and ordering. Classical-mesh dependence is visible in the
solver's own iteration counts across refinements, which are in
\texttt{results/tensor\_verif\_*.csv}.

\end{document}
